\begin{block}{EATS: Efficiency-Aware Task Score}
  A single score that rewards architectures preserving answer quality while
  avoiding unnecessary large-model usage and cost.

  \vspace{4pt}
  \begin{center}
  \begin{tikzpicture}
    \node[draw=posterTeal, fill=posterTeal!12, rounded corners=6pt, inner sep=12pt,
          text width=0.88\linewidth, align=center] {
      \[
        \mathrm{EATS} \;=\;
        \frac{\text{accuracy}}
             {\text{llm\_call\_ratio} \times \text{normalized\_cost} + \varepsilon}
      \]
    };
  \end{tikzpicture}
  \end{center}

  \vspace{2pt}
  \begin{itemize}
    \item \textbf{accuracy} --- fraction of correct answers.
    \item \textbf{llm\_call\_ratio} --- share of queries that trigger an LLM call
          ($0=$ fully SLM, $1=$ always LLM).
    \item \textbf{normalized\_cost} --- total cost $\div$ full-LLM-baseline cost.
    \item $\boldsymbol{\varepsilon}$ --- small constant ($0.01$) avoiding division by zero.
  \end{itemize}
  Higher is better: a system at $0.75$ accuracy that calls the LLM only $10\%$ of
  the time scores far above $0.80$ accuracy that always uses the LLM.

  \vspace{4pt}
  \textbf{Measured per request} via a vLLM metrics proxy (\texttt{\_thesis} block):
  latency, input/output tokens, infra cost, \texttt{energy\_kwh}, and
  \texttt{co2\_g} (grid factor $380$\,g\,CO\textsubscript{2}/kWh by default).
\end{block}
